% =============================================================================
% Kapitel 4: Implementierung
% =============================================================================
\chapter{Implementierung}
\label{ch:implementierung}

\todo{Einleitenden Text verfassen. Was wird in diesem Kapitel vorgestellt?}

% -----------------------------------------------------------------------------

\section{Anforderungen}
%table einfügen für requirements

Im folgenden werden die funktionalen und nicht-funktionalen Anforderungen an das System beschrieben. Diese Anforderungen bilden die Grundlage für die Implementierung und dienen als Referenzpunkt für die Bewertung des Systems. Die Anforderungen werden Anhand von Titel, Kürzel, Kurzbeschreibung und Priorität dargestellt. Die Priorität wird in drei Stufen unterteilt: \textit{hoch}, \textit{mittel} und \textit{niedrig}.

\begin{table}[htbp]
    \centering
    \caption{Anforderungen}
    \label{tab:anforderungen}
    \begin{tabularx}{\textwidth}{|l|l|X|l|}
        \hline
        \textbf{Titel} & \textbf{Kürzel} & \textbf{Kurzbeschreibung} & \textbf{Priorität} \\
        \Xhline{1pt}
        \multicolumn{4}{|l|}{\textbf{Daten}} \\
        \Xhline{1pt}
        Erhebung    & R-1.0 & Daten für die Erstellung von Modellen werden erfasst, auf das richtige Format überprüft und bereitgestellt. & hoch   \\
        \hline
        Aufbereitung & R-1.1 & Daten werden aufbereitet und in das richtige Format gebracht: Informationsgehalt und -qualität, Granularität & mittel \\
        \hline
        Speicherung  & R-1.2 & Daten werden dauerhaft persistiert. & niedrig \\
        \Xhline{1pt}
        \multicolumn{4}{|l|}{\textbf{System}} \\
        \Xhline{1pt}
        Systemstart      & R-2.0 & Funktionalität 4 & hoch   \\
        \hline
        Systemüberwachung & R-2.1 & Funktionalität 5 & mittel \\
        \Xhline{1pt}
        \multicolumn{4}{|l|}{\textbf{Machine Learning Modelle}} \\
        \Xhline{1pt}
        Modelltraining   & R-3.0 & Funktionalität 6 & hoch   \\
        \hline
        Modellevaluierung & R-3.1 & Funktionalität 7 & hoch \\
        \hline
    \end{tabularx}
\end{table}

\section{Systemarchitektur}
\label{sec:systemarchitektur}

\todo{Die Gesamtarchitektur des Systems beschreiben. Welche Komponenten gibt es und wie interagieren sie?}

% -----------------------------------------------------------------------------
\section{Technologieauswahl}
\label{sec:technologieauswahl}

\todo{Welche Technologien, Frameworks, Programmiersprachen und Werkzeuge wurden eingesetzt und warum?}

% -----------------------------------------------------------------------------
\section{Umsetzung}
\label{sec:umsetzung}

\todo{Zentrale Aspekte der Implementierung beschreiben. Wichtige Algorithmen, Datenstrukturen oder Designentscheidungen erläutern.}

% Beispiel für ein Code-Listing:
% \begin{lstlisting}[language=Python, caption={Beispiel-Code}, label={lst:beispiel}]
% def hello_world():
%     print("Hallo Welt!")
% \end{lstlisting}

% -----------------------------------------------------------------------------
\section{Herausforderungen}
\label{sec:herausforderungen}

\todo{Welche Herausforderungen traten bei der Implementierung auf? Wie wurden sie gelöst?}

